\documentclass[12pt,a4paper]{article}
\usepackage[spanish]{babel} % Corta palabras en espa~nol
\usepackage[utf8]{inputenc} % Escribir con acentos, ~n...
\usepackage{eurosym} % smbolo del euro

\usepackage{graphicx}
\begin{document}
%\secttion{Introducción}
\leftline{\textbf{\huge{Ejemplos de Fórmulas Matemáticas y Figuras}}}

\textbf{Veamos ahora las figuras:}

\begin{figure}[h]
\begin{center}
\includegraphics[scale=0.075]{tiger.pdf} \end{center}
\caption{Tigre}
\end{figure}

\textbf{{Elementos flotantes}}
 son tablas y figuras, que pueden estar flotantes en el documento. La posición se especifica:
\begin{itemize}
\item h aquí
\item t Comienzo de una página de texto
\item b Final de una página de texto
\item p En una página sin texto
\end{itemize}
La forma usual es:
\begin{figure}[b]
\begin{center}
\includegraphics[scale=0.075, angle=45]{tiger.pdf} \end{center}
\caption{Flotante Tigre}
\end{figure}


%\medskip
Este primer ejemplo trata de demostrar la facilidad de
\LaTeX{} para hacer \textbf{ecuaciones}:
\[
x=\frac{a_2 x^2 + a_1 x + a_0}{1+2z^3}, \quad
x+y^{2n+2}=\sqrt{b^2-4ac}
\]




\[
\left[ \begin{array}{ccc}
a_{11} & a_{12} & a_{13} \\
a_{21} & a_{22} & a_{23} \\
a_{31} & a_{32} & a_{33}
\end{array} \right]
\left\{ \begin{array}{c}
x_1 \\ x_2 \\ x_3
\end{array} \right\} =
\left( \begin{array}{c}
b_1 \\ b_2 \\ b_3
\end{array} \right)
\]

Para numerar ecuaciones y referenciarlas posteriormente, utilizar los
ambientes:

\begin{eqnarray}
y &=& 1+x+x^2 +x^3 + \cdots \nonumber \\
&=& 1+x\left(1+x+x^2+\cdots \right)
\label{ecu1}
\end{eqnarray}

\end{document}